% 第13章 非线性{$\sigma$}模型：弱耦合
\chapter{非线性 $\sigma$ 模型：弱耦合}

\section{格子正规化}

上一章中，$d$ 维非线性 {$\sigma$} 模型（NLSM）作为 $d-1$ 维量子 Heisenberg 反铁磁体（附加 Berry 相位）的连续近似出现。配分函数为

\begin{equation*}
Z_{\mathrm{NLSM}} = \int_{\Lambda} \mathcal {D} \hat {\mathbf {n}} \, \exp \left( - \frac {\Lambda^ {d - 2}}{2 f} \int d ^ {d} x \sum_ {\mu = 1} ^ {d} \partial_ {\mu} \hat {\mathbf {n}} \cdot \partial_ {\mu} \hat {\mathbf {n}} \right).
\tag{13.1}
\end{equation*}

$|\hat{\mathbf{n}}(\mathbf{x})|=1$ 是单位矢量，$x\in R^{d}$。$\Lambda$ 是动量截断，即 $\mathcal{D}\hat{\mathbf{n}}$ 中包含的最短波长。在指定正规化方案之前，路径积分 (13.1) 没有良定义。考虑 $d$ 维立方格子上的经典 Heisenberg 模型，$\mathcal{N}$ 个格点、晶格常量 $a$。其配分函数为

\begin{equation*}
Z_{\mathrm{HM}} = \int \prod_ {i = 1} ^ {\mathcal {N}} d \hat {\Omega} _ {i} \, \exp \left( \frac {\bar {J}}{T} \sum_{\langle ij \rangle} \hat {\Omega} _ {i} \cdot \hat {\Omega} _ {j} \right).
\tag{13.2}
\end{equation*}

$Z_{\mathrm{HM}}$ 的连续极限由替换

\begin{equation*}
\hat {\Omega} _ {i} \rightarrow \eta_ {i} \hat {\mathbf {n}} ( \mathbf {x} _ {i} ) ,
\tag{13.3}
\end{equation*}

给出，其中对铁磁体（反铁磁体）$\eta_{i}$ 为 $1$（$e^{i\mathbf{x}_{i}\cdot\bar{\pi}}$）。求和与差分替换为

\begin{align*}
\sum_ {i} F ( \mathbf {x} _ {i} ) & \rightarrow a ^ {- d} \int d ^ {d} x \, F ( \mathbf {x} ) ,\\
\hat {\Omega} ( \mathbf {x} _ {i} + a \hat {\boldsymbol {x}} _ {\mu} ) - \hat {\Omega} ( \mathbf {x} _ {i} ) & \rightarrow a \, \partial_ {\mu} \hat {\mathbf {n}} ,
\tag{13.4}
\end{align*}

无量纲耦合常数为

\begin{equation*}
\frac {T}{\overline {J}} \leftrightarrow f.
\tag{13.5}
\end{equation*}

测度替换为

\begin{equation*}
\mathcal {D} \hat {\Omega} \rightarrow \prod_{| \mathbf {q} | \leq \Lambda} d \hat {\mathbf {n}} _ {\mathbf {q}} ,
\tag{13.6}
\end{equation*}

其中 $\Lambda$ 是球型（sph）布里渊区半径，与格子的立方区（cube）有相同的自由度数：

\begin{equation*}
\mathcal {N} ^ {- 1} \sum_{| k _ {\mu} | \leq \frac {\pi}{a}} ^ {\mathrm{cube}} = \mathcal {N} ^ {- 1} \sum_{| k | \leq \Lambda} ^ {\mathrm{sph}} = 1.
\tag{13.7}
\end{equation*}

这定出球型布里渊区半径

\begin{equation*}
\Lambda = \left\{ \begin{array}{ll} \pi / a & d {=} 1 \\ 2 \sqrt {\pi} / a & d {=} 2 \\ ( 6 \pi^ {2} ) ^ {1/3} / a & d {=} 3 \end{array} \right..
\tag{13.8}
\end{equation*}

把 (13.2) 中的哈密顿量按 $\hat{n}$ 的梯度展开并保留到二次项，在低温（$f \ll 1$）下建立关系

\begin{equation*}
Z_{\mathrm{HM}} ( T / \bar {J}, a ) \approx Z_{\mathrm{NLSM}} ( f, \Lambda^ {- 1} ) \, e ^ {d \mathcal {N} / f}
\tag{13.9}
\end{equation*}

因子 $e^{d\mathcal{N}/f}$ 是经典基态的 Boltzmann 权重。只要自旋关联长度 $\xi$ 远大于 $\Lambda^{-1}$，连续近似对生成泛函（见附录 B）与长波自旋关联同样成立。在这一区域，NLSM 的性质弱依赖于正规化方案，因此可作为各种 Heisenberg 模型的模型。这些模型的关键特征是其基态流形的 $O(3)$ 对称性与短程相互作用；其他细节主要影响 $f$ 与 $\Lambda$ 的选取。

NLSM 的“非线性”来自模方约束 $|\hat{\mathbf{n}}(\mathbf{x})|=1$。比如说，可以把第一个分量用其他分量写出：

\begin{equation*}
n ^ {(1)} ( \mathbf {x} ) = \sqrt {1 - \left[ n ^ {(2)} ( \mathbf {x} ) \right] ^ {2} - \left[ n ^ {(3)} ( \mathbf {x} ) \right] ^ {2}}.
\tag{13.10}
\end{equation*}

测度于是替换为

\begin{equation*}
\int \mathcal {D} \hat {\mathbf {n}} = \int \mathcal {D} n ^ {(2)} \, \mathcal {D} n ^ {(3)} \, \frac {\Theta \left[ 1 - ( n ^ {(2)} ) ^ {2} - ( n ^ {(3)} ) ^ {2} \right]}{\sqrt {1 - ( n ^ {(2)} ) ^ {2} - ( n ^ {(3)} ) ^ {2}}}.
\tag{13.11}
\end{equation*}

约束引入场分量之间的非谐相互作用。因此可以预期 NLSM 的关联不同于自由（Gauss）模型——后者的测度是 $\mathcal{D}n^{(2)}\mathcal{D}n^{(3)}$。模方约束的一个重要效应是二维中拓扑稳定激发（“孤子”）的存在\footnote{见 Skyrme，及 Belavin 与 Polyakov，文献注释。}。

作为经典 Heisenberg 模型的连续极限，$O(3)$ NLSM 必须服从 6.2 节证明的 Mermin--Wagner 定理：$d = 1, 2$ 时任何有限温度（即 $f > 0$）下禁止自发破缺对称。

下一节我们将演示，如何用 $f$ 的弱耦合展开看到 $O(n)$ NLSM 被序参量涨落破坏有序。

\section{弱耦合展开}

首先把 NLSM 从 $O(3)$ 推广到 $O(n)$ 对称。由于下面的计算对一般 $n$ 不增加任何额外工作量，不妨保持 $n$ 任意。定义 $n$ 个正交归一基矢量：

\begin{align*}
\hat {\mathbf {e}} ^ {\alpha} \, , \quad & \alpha = 0, \dots , n - 1 ,\\
n ^ {\alpha} & = \hat {\mathbf {n}} \cdot \hat {\mathbf {e}} ^ {\alpha} .
\tag{13.12}
\end{align*}

配分函数为

\begin{align*}
Z_{O(n)} & = \int \mathcal {D} \hat {\mathbf {n}} \, \exp \left[ - \frac {\Lambda^ {d - 2}}{2 f} \int d ^ {d} x \sum_ {\mu \alpha} \left( \partial_ {\mu} n ^ {\alpha} \right) ^ {2} \right] ,\\
\mathcal {D} \hat {\mathbf {n}} & = \prod_ {\alpha} \mathcal {D} n ^ {\alpha} \, \delta \left[ 1 - \sum_ {\alpha = 0} ^ {n - 1} \left( n ^ {\alpha} \right) ^ {2} \right] .
\tag{13.13}
\end{align*}

低温即弱耦合 $f \ll 1$ 时，预期主导组态接近某个有序基态。这要求我们破缺路径积分的 $O(n)$ 对称，选一个特定的均匀基态组态，如

\begin{equation*}
\bar {\hat {\mathbf {n}}} = \hat {\mathbf {e}} ^ {0}.
\tag{13.14}
\end{equation*}

小涨落 $|\hat{n}-\hat{e}^{0}| \ll 1$ 参数化为

\begin{equation*}
\hat {\mathbf {n}} = \hat {\mathbf {e}} _ {0} \sqrt {1 - \bar {\phi} ^ {2}} + \sum_ {a = 1} ^ {n - 1} \phi^ {a} \hat {\mathbf {e}} ^ {a} ,
\tag{13.15}
\end{equation*}

其中

\begin{equation*}
\bar {\phi} ^ {2} \equiv \sum_ {a = 1} ^ {n - 1} \left( \phi^ {a} \right) ^ {2}.
\tag{13.16}
\end{equation*}

这里及以后，拉丁指标（如 $a, b$）表示横向涨落，不包括纵向 $\hat{e}_{0}$ 方向。

NLSM 按 $|\delta\hat{n}|$ 的领头阶展开：

\begin{align*}
Z & \approx \int \mathcal {D} \phi \, \exp \left( - \mathcal {S} ^ {(2)} [ \phi ] \right) \left[ 1 + \mathcal {O} ( f ) \right] ,\\
\mathcal {S} ^ {(2)} & = \frac {\Lambda^ {d - 2}}{2 f} \int d ^ {d} x \sum_ {a, \mu} \left( \partial_ {\mu} \phi^ {a} \right) ^ {2} + \mathcal {O} ( \phi^ {3} )\\
& = \frac {1}{2 f \Lambda^ {2} \mathcal {N}} \sum_{\mathbf {k}, a} \mathbf {k} ^ {2} \phi_{\mathbf {k}} ^ {a} \phi_{- \mathbf {k}} ^ {a} ,
\tag{13.17}
\end{align*}

其中

\begin{equation*}
\phi_ {\mathbf {k}} ^ {a} = \Lambda^ {d} \int d ^ {d} x \, e ^ {- i \mathbf {k} \cdot \mathbf {x}} \phi^ {a} ( \mathbf {x} ).
\tag{13.18}
\end{equation*}

$\mathcal{S}^{(2)}$ 是谐波或“自旋波”能量泛函。用 (13.15) 与 (13.17) 可以计算自旋波对局域涨落的贡献：

\begin{align*}
\langle \left| \delta \hat {\mathbf {n}} ( 0 ) \right| ^ {2} \rangle & \approx \frac {1}{Z \mathcal {N}} \sum_{\mathbf {k} a} \int \mathcal {D} \hat {\mathbf {n}} \, \phi_{\mathbf {k}} ^ {a} \phi_{- \mathbf {k}} ^ {a} \exp \left( - \mathcal {S} ^ {(2)} [ \phi ] \right)\\
& = \frac {( n - 1 ) \Lambda^ {2} f}{( 2 \pi ) ^ {d}} \lim_{\tilde {\Lambda} \to 0} \int_{\tilde {\Lambda}} ^ {\Lambda} \frac {d ^ {d} k}{| \mathbf {k} | ^ {2}}\\
& \sim \left\{ \begin{array}{ll} ( n - 1 ) f \lim_{\tilde {\Lambda} \to 0} \tilde {\Lambda} ^ {- 1} = \infty & d = 1 \\ - \dfrac {( n - 1 ) f}{2 \pi} \lim_{\tilde {\Lambda} \to 0} \ln \tilde {\Lambda} = \infty & d = 2 \\[4pt] \dfrac {( n - 1 ) f \Lambda}{2 \pi^ {2}} < \infty & d = 3 \end{array} \right.
\tag{13.19}
\end{align*}

$n \geq 2$ 时，一维、二维的自旋波涨落对所有 $f > 0$ 都发散。于是朴素的弱耦合展开失败——自发破缺对称的假设不成立。这一点我们凭 Mermin--Wagner 定理 6.2 早已知道。由 Haldane 映射，一维量子 Heisenberg 反铁磁体的基态对所有有限自旋大小 $S < \infty$ 都应是无序的（自旋液体）。

我们用“朴素展开”一词，暗示更精细的方法可以处理 (13.19) 的红外发散——见下节。

\section{穷人重整化}

(13.19) 的红外发散表明，必须放弃破缺对称的假设，保留路径积分的 $O(n)$ 对称。但仍然可以利用 $f$ 的小来近似无序相的关联。做法是把被积函数绕缓变组态 $\hat{\mathbf{n}}^{0}(\mathbf{x})$ 展开（后者尚待积分）。逐次消去大动量模式（“快模式”），重整化剩余“慢”模式的能量泛函。本节推导 $\beta$ 函数所体现的重整化群变换；下节用 $\beta$ 函数确定关联长度。

循 Polyakov，把 $\hat{n}$ 分成慢自由度 $\hat{n}^{0}$ 与快自由度 $\phi^{a}$：

\begin{align*}
\hat {\mathbf {n}} ( \mathbf {x} ) & = \hat {\mathbf {n}} ^ {0} ( \mathbf {x} ) \sqrt {1 - \bar {\phi} ^ {2}} + \sum_ {a = 1} ^ {n - 1} \phi^ {a} \hat {\mathbf {e}} ^ {a} ( \mathbf {x} ) ,\\
\bar {\phi} ^ {2} & = \sum_ {a = 1} ^ {n - 1} \left( \phi^ {a} \right) ^ {2} ,
\tag{13.20}
\end{align*}

其中随动基底定义为

\begin{align*}
\hat {\mathbf {e}} ^ {0} ( \mathbf {x} ) & \equiv \hat {\mathbf {n}} ^ {0} ( \mathbf {x} ) ,\\
\hat {\mathbf {e}} ^ {\alpha} ( \mathbf {x} ) \cdot \hat {\mathbf {e}} ^ {\beta} ( \mathbf {x} ) & = \delta_ {\alpha \beta} \, , \quad \alpha = 0, 1, \ldots , n - 1 .
\tag{13.21}
\end{align*}

希腊指标包括纵向方向 $\hat{n}^{0}$；拉丁指标限于横向 $a, b = 1, \ldots, n - 1$。

空间相邻点坐标系之间的关系定义规范势\footnote{在微分几何中称为联络。}

\begin{equation*}
\tilde {A} _ {\mu} ^ {\alpha \beta} \equiv \hat {\mathbf {e}} ^ {\alpha} \cdot \partial_ {\mu} \hat {\mathbf {e}} ^ {\beta} ,
\tag{13.22}
\end{equation*}

它描述活动正交归一基的变化：

\begin{align*}
\partial_ {\mu} \hat {\mathbf {n}} ^ {0} & = \sum_ {a} \tilde {A} _ {\mu} ^ {a 0} \hat {\mathbf {e}} ^ {a} ,\\
\partial_ {\mu} \hat {\mathbf {e}} ^ {a} & = \sum_ {b} \left( \tilde {A} _ {\mu} ^ {ba} \hat {\mathbf {e}} ^ {b} \right) - \tilde {A} _ {\mu} ^ {a 0} \hat {\mathbf {n}} ^ {0} .
\tag{13.23}
\end{align*}

由于 $\hat{\mathbf{e}}^{\alpha}$ 正交归一，$\tilde{A}_{\mu}^{\alpha \beta}$ 必须反对称：

\begin{equation*}
\tilde {A} _ {\mu} ^ {\alpha \beta} = - \tilde {A} _ {\mu} ^ {\beta \alpha}.
\tag{13.24}
\end{equation*}

记一个马上要用的恒等式：

\begin{equation*}
\sum_ {a} \left( \tilde {A} _ {\mu} ^ {a 0} \right) ^ {2} = \left( \partial_ {\mu} \hat {\mathbf {n}} ^ {0} \right) ^ {2}.
\tag{13.25}
\end{equation*}

选中间动量尺度 $\tilde{\Lambda}$：

\begin{equation*}
0 < \tilde {\Lambda} \ll \Lambda ,
\tag{13.26}
\end{equation*}

区分“快”场

\begin{equation*}
\phi^ {a} ( \mathbf {x} ) \equiv \sum_{\tilde {\Lambda} \leq | \mathbf {k} | \leq \Lambda} \exp [ i \mathbf {kx} ] \, \phi_{\mathbf {k}} ^ {a} ,
\tag{13.27}
\end{equation*}

与慢场

\begin{equation*}
\mathbf {X} ( \mathbf {x} ) = \sum_{0 \leq | \mathbf {k} | < \tilde {\Lambda}} \exp [ i \mathbf {kx} ] \, \mathbf {X} _ {\mathbf {k}} \, , \quad \mathbf {X} = \hat {\mathbf {e}} ^ {\alpha}, \tilde {A} _ {\mu} ^ {\alpha , \beta}.
\tag{13.28}
\end{equation*}

把 (13.20) 代入 NLSM Lagrange 量：

\begin{align*}
\mathcal {L} & = \frac {\Lambda^ {d - 2}}{2 f} \int d ^ {d} x \sum_{\mu = 1} ^ {d} \partial_ {\mu} \hat {\mathbf {n}} \cdot \partial_ {\mu} \hat {\mathbf {n}}\\
& = \frac {\Lambda^ {d - 2}}{2 f} \sum_{\mu ab} \Big[ \left( \partial_ {\mu} \sqrt {1 - \bar {\phi} ^ {2}} - \tilde {A} _ {\mu} ^ {a 0} \phi^ {a} \right) ^ {2}\\
& \quad + \left( \partial_ {\mu} \phi^ {a} + \tilde {A} _ {\mu} ^ {ab} \phi^ {b} + \tilde {A} _ {\mu} ^ {a 0} \sqrt {1 - \bar {\phi} ^ {2}} \right) ^ {2} \Big] .
\tag{13.29}
\end{align*}

把 Lagrange 量\footnote{变换 (13.20) 的 Jacobian 不贡献二次阶修正。}按 $\phi$ 展开到二次阶：

\begin{align*}
Z & \approx \int_{\tilde {\Lambda}} \mathcal {D} \hat {\mathbf {n}} ^ {0} \exp \left( - \int_{\tilde {\Lambda}} d ^ {d} x \, \mathcal {L} [ \hat {\mathbf {n}} ^ {0} ] \right) Z ^ {(2)} [ \hat {\mathbf {n}} ^ {0} ] \left[ 1 + \mathcal {O} ( \phi^ {2}, f ^ {- 1} \phi^ {3} ) \right] ,\\
Z ^ {(2)} & = \int_{\Lambda} \mathcal {D} \phi \exp \left( - \int_{\Lambda} d ^ {d} x \, \mathcal {L} ^ {(2)} \left[ \hat {\mathbf {n}} ^ {0}, \phi \right] \right) ,
\tag{13.30}
\end{align*}

关系 (13.25) 给出 $L[\hat{n}^{0}]$ 项，快（“自旋波”）Lagrange 量为

\begin{equation*}
\mathcal {L} ^ {(2)} = \frac {\Lambda^ {d - 2}}{2 f} \left[ \sum_{\mu ab} \left( \partial_ {\mu} \phi^ {a} + \tilde {A} _ {\mu} ^ {ab} \phi^ {b} \right) ^ {2} + \tilde {A} _ {\mu} ^ {a 0} \tilde {A} _ {\mu} ^ {b 0} \left( \phi^ {a} \phi^ {b} - \delta_ {ab} \bar {\phi} ^ {2} \right) \right].
\tag{13.31}
\end{equation*}

我们忽略 $\phi$ 的三次及更高次幂，把分析限于 $f$ 的领头阶项。用重整化群的行话，我们在做“单圈水平”的计算。

在 (13.31) 中我们还忽略了耦合快慢场的交叉项

\begin{equation*}
\tilde {A} _ {\mathbf {q}} ^ {ab} \phi_{\mathbf {k - q}} ^ {b} \tilde {A} _ {- \mathbf {k}} ^ {a 0} ,
\tag{13.32}
\end{equation*}

选 $\tilde{\Lambda} \ll \Lambda$ 保证这类项的贡献很小：它们对所有 $k > 2\tilde{\Lambda}$（即布里渊区的大部分）为零。

$\mathcal{L}^{(2)}$ 描述绕 $\hat{\mathbf{n}}^{0}$ 的高动量涨落。警觉的读者会注意到，$n > 2$ 时 (13.20) 中横向单位矢量 $\hat{\mathbf{e}}^{a}$ 的选取有某种任意性：可在空间每点用任何正交 $(n - 1) \times (n - 1)$ 矩阵 $R$ 连续转动 $n - 1$ 维子空间中的横向坐标系：

\begin{equation*}
\phi^ {a} \rightarrow R ^ {ab} \phi^ {b} \, , \quad R ^ {T} R = 1 ,
\tag{13.33}
\end{equation*}

它把 (13.31) 第一项中的协变导数变换为

\begin{equation*}
\partial_ {\mu} \delta_ {ab} + \tilde {A} _ {\mu} ^ {ab} \rightarrow \partial_ {\mu} \delta_ {ab} + \tilde {A} _ {\mu} ^ {ab} + \left[ ( \partial_ {\mu} R ) R ^ {- 1} \right] ^ {ab}.
\tag{13.34}
\end{equation*}

$\mathcal{L}^{(2)}$ 在 (13.34) 下的不变性是快模式 Lagrange 量的一个规范对称。容易看出 $\mathcal{L}^{(2)}$ 只能依赖 $\tilde{A}_{\mu}^{ab}$ 的导数；由于这些势是慢场，这些项是 $\hat{n}^{0}$ 的比领头项 $L[\hat{n}^{0}]$ 更高阶的导数。因此下面的讨论忽略 $\mathcal{L}^{(2)}$ 中的规范势。习题中将看到，当绕具有有限拓扑密度的慢场展开时，规范势是重要的。

在 (13.30) 中完成 Gauss 积分：

\begin{equation*}
Z ^ {(2)} [ \hat {\mathbf {n}} ^ {0} ] \propto \exp \left[ - \frac {1}{2} \mathrm{Tr} \ln \left( \Pi_ {0} - \Pi_ {1} \right) \right].
\tag{13.35}
\end{equation*}

整体归一化常数与 $\hat{n}^{0}$ 无关，对关联函数不重要。$\Pi_{0}$ 与 $\Pi_{1}$ 是快模式空间 $|k| > \tilde{\Lambda}$ 中的矩阵：

\begin{align*}
\left( \Pi_ {0} \right)_{\mathbf {k}, \mathbf {k} '} & = - \left( \partial_ {\mu} ^ {2} \right) \delta_ {ab} = k ^ {2} \delta_{\mathbf {k}, - \mathbf {k} '} \delta_ {ab} ,\\
\left( - \Pi_ {1} \right)_{\mathbf {k}, \mathbf {k} '} & = \delta_{\mathbf {k}, - \mathbf {k} '} \left[ \tilde {A} _ {\mu} ^ {a 0} \tilde {A} _ {\mu} ^ {b 0} - \delta_ {ab} \sum_ {c} \left( \tilde {A} _ {\mu} ^ {c 0} \right) ^ {2} \right] .
\tag{13.36}
\end{align*}

用恒等式

\begin{equation*}
- \frac {1}{2} \mathrm{Tr} \ln \left( \Pi_ {0} - \Pi_ {1} \right) = - \frac {1}{2} \mathrm{Tr} \ln \Pi_ {0} + \frac {1}{2} \sum_ {n = 1} ^ {\infty} \frac {1}{n} \mathrm{Tr} \left( \Pi_ {0} ^ {- 1} \Pi_ {1} \right) ^ {n}.
\tag{13.37}
\end{equation*}

把 $Z^{(2)}$ 中的指数按 $\tilde{A}^{a0}$ 幂次展开。与先前的近似一致，忽略更高阶导数，只保留 $n = 1$ 项：

\begin{align*}
\frac {1}{2} \mathrm{Tr} \left( \Pi_ {0} ^ {- 1} \Pi_ {1} \right) & = - \frac {1}{2 \mathcal {N}} \sum_{\tilde {\Lambda} < \mathbf {k} < \Lambda} \frac {1}{k ^ {2}} \sum_ {\mu a} \left[ \left( \tilde {A} _ {\mu} ^ {a 0} \right) ^ {2} - ( n - 1 ) \left( \tilde {A} _ {\mu} ^ {a 0} \right) ^ {2} \right]\\
& \approx \frac {\Lambda^ {d - 2} \Delta_ {d}}{2} \sum_ {\mu} \left( \partial_ {\mu} \hat {\mathbf {n}} ^ {0} \right) ^ {2} ,
\tag{13.38}
\end{align*}

这里用了 (13.25)，并定义无量纲常数

\begin{align*}
\Delta_ {d} ( \tilde {\Lambda} / \Lambda ) & = ( n - 2 ) \Lambda^ {2 - d} \int_{\tilde {\Lambda} < | \mathbf {k} | < \Lambda} \frac {d ^ {d} k}{( 2 \pi ) ^ {d}} \frac {1}{k ^ {2}}\\
& = \left\{ \begin{array}{ll} \dfrac {( n - 2 )}{\zeta_ {1}} \left[ ( \tilde {\Lambda} / \Lambda ) ^ {d - 2} - 1 \right] & d = 1 \\[4pt] - \dfrac {n - 2}{\zeta_ {2}} \ln ( \tilde {\Lambda} / \Lambda ) & d = 2 \\[4pt] \dfrac {( n - 2 )}{\zeta_ {3}} \left[ 1 - ( \tilde {\Lambda} / \Lambda ) ^ {d - 2} \right] & d = 3 \end{array} \right. ,
\tag{13.39}
\end{align*}

其中 $\zeta_1 = \pi$（原书如此），$\zeta_2 = 2\pi$，$\zeta_3 = 2\pi^2$。

由 (13.38)，我们发现积掉快模式对 NLSM 的修正正比于 NLSM Lagrange 量本身！把修正并入 NLSM 并替换截断 $\Lambda \rightarrow \tilde{\Lambda}$：

\begin{equation*}
Z \approx \int_{\tilde {\Lambda}} \mathcal {D} \hat {\mathbf {n}} ^ {0} \exp \left[ - \frac {\tilde {\Lambda} ^ {d - 2}}{2 \tilde {f}} \int_{\tilde {\Lambda}} d ^ {d} x \left( \partial_ {\mu} \hat {\mathbf {n}} ^ {0} \right) ^ {2} \right] ,
\tag{13.40}
\end{equation*}

重整化耦合常数为

\begin{equation*}
\tilde {f} = ( \tilde {\Lambda} / \Lambda ) ^ {d - 2} \frac {f}{1 - f \Delta_ {d} ( \tilde {\Lambda} / \Lambda )}.
\tag{13.41}
\end{equation*}

方程 (13.41) 在深入的微扰区域 $f, \tilde{f} \ll 1$ 有效。否则，像 (13.30) 中那样忽略 Lagrange 量的三次及更高项就不合理了。

$d > 2$ 时，$\Delta$ 在红外极限 $\tilde{\Lambda} \to 0$ 不发散，$\tilde{f}$ 的修正很小。反之，$d \leq 2$ 且 $n \geq 3$ 时，$\Delta_d$ 在红外极限发散。

比较 (13.19) 与 (13.39) 很有启发：虽然 $d=1,2$ 时基态涨落按 $n-1$ 发散，(13.39) 中 $\Delta_{d}$ 的红外发散只在 $n\geq3$ 时发生。可见 $n=2$ 情形特殊：只有一个横向自旋波模式，单圈水平上自旋波理论实际上无相互作用；而 $n>2$ 时自旋波模式间的相互作用在重整化下流向强耦合区域。

$O(2)$ NLSM 是 $xy$ 模型的连续极限：

\begin{equation*}
\mathcal {H}_{xy} = \bar {J} \sum_{\langle ij \rangle} \left( \hat {\Omega} _ {i} ^ {x} \hat {\Omega} _ {j} ^ {x} + \hat {\Omega} _ {i} ^ {y} \hat {\Omega} _ {j} ^ {y} \right) \sim \mathcal {L}_{O ( 2 )}.
\tag{13.42}
\end{equation*}

$d=2$ 时 $\mathcal{L}_{O(2)}$ 的 $\beta$ 函数由 $f$ 的更高阶（三次）幂主导，上面未计算。结果是：连续理论存在弱耦合相，$f>0$ 有“准长程序”，即按幂律衰减的关联。更高温度（更大 $f$）时连续近似 (13.42) 失效，涡旋奇性变得重要。Kosterlitz 与 Thouless（KT）证明：涡旋对的解绑在 Kosterlitz--Thouless 温度处摧毁准长程序，此后自旋关联在长距离按指数衰减。$O(2)$ 模型与 KT 相变在许多物理现象（包括超流与超导）中极为重要，文献中有详尽综述。这里我们把进一步的讨论限于 $n \geq 3$ 的 Heisenberg 型模型。

总之，到 $f$ 的领头阶，在改变上截断时 NLSM 标度到自身，我们（近似地）得到了耦合常数的重整化方程。我们的分析假设重整化时唯一增长的耦合常数是 $f$，即更高阶相互作用与更高阶导数项要么“无关”要么“边缘”。这一假设应由显式计算验证；这里我们欣然略过这一问题——这正是这类分析被称为“穷人标度”的原因\footnote{该词由 P. W. Anderson 创造，意指 Anderson、Yuval 与 Hamman 对 Kondo 模型的重整化。}。用 $d-2$ 展开对生成泛函更仔细的处理由 Brézin 与 Zinn-Justin 完成。

\section{$\beta$ 函数}

这里去掉 $\tilde{f}$ 上的波浪号，把依赖截断的耦合常数记作 $f(\Lambda)$。(13.41) 给出了 $\Lambda \rightarrow \tilde{\Lambda}$ 时 $f$ 的领头阶重整化。支配耦合常数在重整化下“流动”的 $\beta$ 函数定义为

\begin{equation*}
\frac {d f}{d \ln ( \tilde {\Lambda} / \Lambda )} \equiv \beta ,
\tag{13.43}
\end{equation*}

由 (13.41) 得

\begin{equation*}
\beta = ( d - 2 ) f - \frac {n - 2}{\zeta_ {d}} f ^ {2}.
\tag{13.44}
\end{equation*}

于是 $\beta$ 显式地只依赖 $f$，不依赖 $\Lambda$ 或 $\tilde{\Lambda}$。这一性质称为“普适性”。$\beta$ 函数的零点是重整化变换的“不动点”：

\begin{equation*}
\beta \left( f _ {c} \right) = 0.
\tag{13.45}
\end{equation*}

\begin{figure}[H]
\centering
\includegraphics[width=0.55\textwidth]{fig13_1.jpg}
\caption{$\beta$ 函数与 $\Lambda$ 减小时 $f$ 的流动。}
\label{fig:13.1}
\end{figure}

$\bar{\Lambda}$ 向下重整化时，$\beta < 0$ 耦合常数减小，$\beta > 0$ 增大。不动点为

\begin{equation*}
\begin{array}{ll}
d = 1 \; : & f _ {c} = 0 \quad \text{不稳定} \, ,\\
d = 2 \; : & f _ {c} = 0 \quad \text{不稳定} \, ,\\
d = 3 \; : & \left\{ \begin{array}{ll} f _ {c} = 0 & \text{稳定} \\ f _ {c} = \dfrac {2 \pi^ {2}}{n - 2} & \text{不稳定} \end{array} \right. .
\end{array}
\tag{13.46}
\end{equation*}

“稳定”与“不稳定”指 $f \approx f_{c}$ 在重整化下的流动趋向或背离不动点。如图~\ref{fig:13.1} 所示，一维、二维中耦合常数流向强耦合；三维存在弱耦合相与强耦合相。$f < f_{c}$ 时耦合流向弱相互作用不动点——微扰论有效，绕有序基态的低阶自旋波理论成立。如图~\ref{fig:13.1}，$f > f_{c}$ 时体系流向强耦合，即无序。$f_{c}$ 是区分破缺对称相与无序相的临界耦合常数。

深入强耦合相，假设依赖动量的结构因子是光滑的，短程序表现为 $q = 0$ 处的极大：

\begin{equation*}
\frac {1}{\mathcal {N}} \langle \hat {\mathbf {n}} _ {\mathbf {q}} \hat {\mathbf {n}} _ {- \mathbf {q}} \rangle \approx \frac {C}{| \mathbf {q} | ^ {2} + \xi^ {- 2} + \mathcal {O} ( | \mathbf {q} ^ {3} | )} ,
\tag{13.47}
\end{equation*}

$C$ 是有量纲常数。(13.47) 的 Fourier 变换是“Ornstein--Zernicke”关联函数，其在大 $|x|/\xi$ 的渐近形式为\footnote{见 Fisher 的论文，文献注释。}

\begin{equation*}
\langle \hat {\mathbf {n}} ( 0 ) \hat {\mathbf {n}} ( \mathbf {x} ) \rangle \propto | \mathbf {x} | ^ {- ( d - 1 )/2} \exp ( - | \mathbf {x} | / \xi ) \left( 1 + \frac {\xi ( d - 3 )}{| \mathbf {x} |} \right) ,
\tag{13.48}
\end{equation*}

关联长度 $\xi$ 支配关联在远距离的指数衰减，是关联的特征长度尺度。用高截断 $\Lambda$ 与裸耦合 $f(\Lambda)$ 的 NLSM 估算 $\xi$，还是用 $\tilde{\Lambda}$ 与 $f(\tilde{\Lambda})$ 的重整化模型估算，本不应有区别。这一不变性可表述为泛函恒等式：

\begin{equation*}
\xi [ f ( \Lambda ), \Lambda ] = \xi [ f ( \tilde {\Lambda} ), \tilde {\Lambda} ].
\tag{13.49}
\end{equation*}

左端对 $\ln (\tilde{\Lambda})$ 的导数为零，故

\begin{equation*}
\frac {\partial \xi}{\partial \ln ( \tilde {\Lambda} )} + \beta \frac {\partial \xi}{\partial f} = 0.
\tag{13.50}
\end{equation*}

由于 $\xi$ 有距离量纲而 $f$ 无量纲，由量纲分析：

\begin{equation*}
\xi ( \tilde {\Lambda}, f ) = \tilde {\Lambda} ^ {- 1} \phi ( f ) ,
\tag{13.51}
\end{equation*}

$\phi$ 是无量纲函数。于是

\begin{equation*}
\frac {\partial \xi}{\partial \ln ( \tilde {\Lambda} )} = - \xi ,
\tag{13.52}
\end{equation*}

从而 (13.50) 有显式解

\begin{align*}
\int_{\xi_ {0}} ^ {\xi} \frac {d \xi}{\xi} & = \int_{f_ {0}} ^ {f} \frac {d f}{\beta ( f )}\\
\rightarrow \xi & = \xi_ {0} \exp \left( \int_{f_ {0}} ^ {f} \frac {d f}{\beta ( f )} \right) .
\tag{13.53}
\end{align*}

$\xi_{0}$ 与 $f_{0}$ 是积分常数，可由强耦合处的关联长度定出（见稍后讨论）。

三维中 $\beta$ 在 $f_{c}$ 处线性消失，$\xi$ 按幂律发散：

\begin{align*}
\xi^ {d = 3} & \sim \left( \frac {f - f _ {c}}{f _ {0}} \right) ^ {1 / \beta ' ( f _ {c} )} ,\\
\beta ' ( f _ {c} ) & = - 1 + \mathcal {O} ( 1 ) .
\tag{13.54}
\end{align*}

二维中由 (13.44)，$\beta$ 函数在 $f_{c}=0$ 处二次消失。可以从强耦合区（$f_{0} \gg 1$，关联长度 $\xi_{0}=O(a)$）到物理 $f$ 值积分 (13.53) 右端：

\begin{equation*}
\xi^ {d = 2} \sim \xi_ {0} \exp \left( \frac {2 \pi}{( n - 2 ) f} \right).
\tag{13.55}
\end{equation*}

对给定格子模型确定 $\xi_{0}$，可以把 (13.55) 与强耦合（高温）下数值定出的关联长度拟合。$\xi_{0}/a$ 是依赖格子模型短程细节的数值常数\footnote{见 Shenker 与 Tobochnik。}。

一维中 (13.53) 积分给出

\begin{equation*}
\xi^ {d = 1} \sim \frac {\xi_ {0} f _ {0}}{f} ,
\tag{13.56}
\end{equation*}

$\xi_{0}f_{0}$ 是 $\Lambda^{-1}$ 的量级。

\begin{exercises}

\item 对 $O(3)$ 模型，把两个自旋波坐标写成一个复场：

\begin{equation*}
\psi ( \mathbf {x} ) = \phi^ {1} ( \mathbf {x} ) + i \phi^ {2} ( \mathbf {x} ) ,
\tag{13.57}
\end{equation*}

并定义 $U(1)$ 规范场

\begin{equation*}
A _ {\mu} [ \hat {\mathbf {n}} ^ {0} ] = - \frac {i}{2} \tilde {A} ^ {12} [ \hat {\mathbf {n}} ^ {0} ].
\tag{13.58}
\end{equation*}

通过对角化 (13.31) 的 $\mathcal{L}^{(2)}$ 中的二次型，求有限规范势存在时的自旋波本征值 $\epsilon_{n}$：

\begin{equation*}
\int d ^ {d} x \left( \partial_ {\mu} \phi^ {a} + \tilde {A} _ {\mu} ^ {ab} \phi^ {b} \right) ^ {2}.
\tag{13.59}
\end{equation*}

证明 $\epsilon_{n}$ 可由方程

\begin{equation*}
\left( \partial_ {\mu} - i 2 A _ {\mu} \right) ^ {2} \psi_ {n} = \epsilon_ {n} \psi_ {n}.
\tag{13.60}
\end{equation*}

得到。注：(13.60) 描述电磁矢量势 $A_{\mu}$ 存在时电荷为 2 的自由玻色子。

\item 证明对 (13.58) 定义的 $O(3)$ NLSM 规范场 $A(\hat{\mathbf{n}})$：

\begin{equation*}
\left( \partial_ {\mu} A _ {\nu} - \partial_ {\nu} A _ {\mu} \right) = \frac {1}{2} \partial_ {\mu} \hat {\mathbf {n}} \times \partial_ {\nu} \hat {\mathbf {n}} \cdot \hat {\mathbf {n}}
\tag{13.61}
\end{equation*}

\item 对二维体系，假设均匀拓扑密度 $B = \partial_{x} \hat{n} \times \partial_{y} \hat{n} \cdot \hat{n}$（如大 Skyrmion 组态中心附近）。选一方便的规范表达 $A_{x}, A_{y}$，并解 (13.60) 的本征值。提示：用均匀磁场中粒子的 Landau 能级理论\footnote{见 L. D. Landau and E. M. Lifshits, \textit{Quantum Mechanics} (Pergamon Press), p. 457，或见 Auerbach, Larson, and Murthy。}

\end{exercises}

\begin{chapbib}
\item 穷人重整化循 Polyakov 的原始推导：
  A. M. Polyakov, Phys. Lett. B \textbf{59}, 79 (1975)；
  A. M. Polyakov, \textit{Gauge Fields and Strings} (Harwood, 1987), 第二章。
\item 非线性 {$\sigma$} 模型到双圈阶的 $d-2$ 展开计算于
  E. Brézin and J. Zinn-Justin, Phys. Rev. B \textbf{14}, 3110 (1976)。
\item 非线性 {$\sigma$} 模型关联长度的有限格子 Monte Carlo 模拟：
  S. H. Shenker and J. Tobochnik, Phys. Rev. B \textbf{22}, 4462 (1980)。
\item Ornstein--Zernicke 关联函数的 Fourier 变换计算于
  M. E. Fisher, Physica \textbf{28}, 172 (1962)。
\item 二维 NLSM 中的孤子（Skyrmion）由
  T. Skyrme, Proc. R. Soc. London Ser. A \textbf{260}, 127 (1961)；
  A. A. Belavin and A. M. Polyakov, JETP Lett. \textbf{22}, 245 (1975)
  发现。
\item 规范势 $A_{ab}^{\mu}$ 对自旋波谱的影响见
  A. Auerbach, B. Larson, and G. Murthy, Phys. Rev. B \textbf{43}, 11515 (1991)。
\item 二维 $xy$ 模型与涡旋效应：
  J. M. Kosterlitz and D. J. Thouless, J. Phys. C \textbf{6}, 1181 (1973)；
  V. L. Berezinskii, Sov. Phys. JETP \textbf{34}, 610 (1972)。
\end{chapbib}
